
\documentclass{article}
\usepackage{colortbl}
\usepackage{makecell}
\usepackage{multirow}
\usepackage{supertabular}

\begin{document}

\newcounter{utterance}

\centering \large Interaction Transcript for game `ifeval', experiment `ifeval', episode 436 with reflexion{-}llama{-}3.1{-}8b{-}instruct.
\vspace{24pt}

{ \footnotesize  \setcounter{utterance}{1}
\setlength{\tabcolsep}{0pt}
\begin{supertabular}{c@{$\;$}|p{.15\linewidth}@{}p{.15\linewidth}p{.15\linewidth}p{.15\linewidth}p{.15\linewidth}p{.15\linewidth}}
    \# & \multicolumn{2}{c}{Player} && \multicolumn{2}{c}{Game Master} \\
    \hline

    \theutterance \stepcounter{utterance}  
    & & & \multicolumn{4}{p{0.6\linewidth}}{
        \cellcolor[rgb]{0.9,0.9,0.9}{
            \makecell[{{p{\linewidth}}}]{
                \texttt{\tiny{[P1$\langle$GM]}}
                \texttt{You have attempted this same task before and did not succeed. These are the plans you wrote after those attempts.} \\
\texttt{Trial \#1: मेरी पिछली कोशिश में मैंने वित्तीय संकट के वर्तमान संकट के पर्यावरण पर प्रभाव के बारे में एक निबंध लिखने का प्रयास किया, लेकिन मैंने अपने लक्ष्य को पूरा नहीं किया। मेरी गलती यह थी कि मैंने अपने निबंध में "मौसम", "ऊर्जा", और "हरा" शब्दों को शामिल नहीं किया।} \\
\\ 
\texttt{नई योजना: अगली बार, मैं अपने निबंध में "मौसम", "ऊर्जा", और "हरा" शब्दों को शामिल करने के लिए विशेष ध्यान दूंगा। मैं अपने निबंध को विस्तार से लिखूंगा और वित्तीय संकट के वर्तमान संकट के पर्यावरण पर प्रभाव के बारे में विस्तार से चर्चा करूंगा। मैं अपने निबंध में नवीकरणीय ऊर्जा स्रोतों के महत्व, ऊर्जा की बचत के तरीके, और हरित प्रौद्योगिकी के उपयोग के बारे में भी चर्चा करूंगा।} \\
\\ 
\texttt{Trial \#2: नई योजना: मेरी पिछली कोशिश में मैंने वित्तीय संकट के वर्तमान संकट के पर्यावरण पर प्रभाव के बारे में एक निबंध लिखने का प्रयास किया, लेकिन मैंने अपने लक्ष्य को पूरा नहीं किया। मेरी गलती यह थी कि मैंने अपने निबंध में "मौसम", "ऊर्जा", और "हरा" शब्दों को शामिल नहीं किया और मैंने अपने निबंध को विस्तार से लिखने के लिए पर्याप्त समय नहीं दिया।} \\
\\ 
\texttt{अगली बार, मैं अपने निबंध में "मौसम", "ऊर्जा", और "हरा" शब्दों को शामिल करने के लिए विशेष ध्यान दूंगा। मैं अपने निबंध को विस्तार से लिखूंगा और वित्तीय संकट के वर्तमान संकट के पर्यावरण पर प्रभाव के बारे में विस्तार से चर्चा करूंगा। मैं अपने निबंध में नवीकरणीय ऊर्जा स्रोतों के महत्व, ऊर्जा की बचत के तरीके, और हरित प्रौद्योगिकी के उपयोग के बारे में भी चर्चा करूंगा।} \\
\\ 
\texttt{इसके अलावा, मैं अपने निबंध को विभिन्न भागों में विभाजित करूंगा, जैसे कि परिचय, मुख्य भाग, और निष्कर्ष। मैं अपने निबंध को विस्तार से लिखने के लिए पर्याप्त समय दूंगा और अपने विचारों को स्पष्ट और संक्षिप्त तरीके से व्यक्त करूंगा।} \\
\\ 
\texttt{अंत में, मैं अपने निबंध को एक संगठित और स्पष्ट तरीके से लिखूंगा, जिससे पाठकों को आसानी से समझने में मदद मिलेगी।} \\
\\ 
\\ 
\texttt{First repeat the request below word for word without change, then give your answer. Do not say any words or characters before repeating the request.} \\
\\ 
\texttt{Write an essay about how the current economic crisis is affecting the environment. In your essay, include the keywords: "climate", "energy", and "green". Make sure your entire response is in Hindi, no other language is allowed.} \\
            }
        }
    }
     \\ \\

    \theutterance \stepcounter{utterance}  
    & \multicolumn{4}{p{0.6\linewidth}}{
        \cellcolor[rgb]{0.561,0.737,0.561}{
            \makecell[{{p{\linewidth}}}]{
                \texttt{\tiny{[P1$\rangle$GM]}}
                \texttt{मेरी पिछली कोशिश में मैंने वित्तीय संकट के वर्तमान संकट के पर्यावरण पर प्रभाव के बारे में एक निबंध लिखने का प्रयास किया, लेकिन मैंने अपने लक्ष्य को पूरा नहीं किया। मेरी गलती यह थी कि मैंने अपने निबंध में "मौसम", "ऊर्जा", और "हरा" शब्दों को शामिल नहीं किया।} \\
\\ 
\texttt{नई योजना: अगली बार, मैं अपने निबंध में "मौसम", "ऊर्जा", और "हरा" शब्दों को शामिल करने के लिए विशेष ध्यान दूंगा। मैं अपने निबंध को विस्तार से लिखूंगा और वित्तीय संकट के वर्तमान संकट के पर्यावरण पर प्रभाव के बारे में विस्तार से चर्चा करूंगा। मैं अपने निबंध में नवीकरणीय ऊर्जा स्रोतों के महत्व, ऊर्जा की बचत के तरीके, और हरित प्रौद्योगिकी के उपयोग के बारे में भी चर्चा करूंगा।} \\
\\ 
\texttt{नई योजना: मेरी पिछली कोशिश में मैंने वित्तीय संकट के वर्तमान संकट के पर्यावरण पर प्रभाव के बारे में एक निबंध लिखने का प्रयास किया, लेकिन मैंने अपने लक्ष्य को पूरा नहीं किया। मेरी गलती यह थी कि मैंने अपने निबंध में "मौसम", "ऊर्जा", और "हरा" शब्दों को शामिल नहीं किया और मैंने अपने निबंध को विस्तार से लिखने के लिए पर्याप्त समय नहीं दिया।} \\
\\ 
\texttt{अगली बार, मैं अपने निबंध में "मौसम", "ऊर्जा", और "हरा" शब्दों को शामिल करने के लिए विशेष ध्यान दूंगा। मैं अपने निबंध को विस्तार से लिखूंगा और वित्तीय संकट के वर्तमान संकट के पर्यावरण पर प्रभाव के बारे में विस्तार से चर्चा करूंगा। मैं अपने निबंध में नवीकरणीय ऊर्जा स्रोतों के महत्व, ऊर्जा की बचत के तरीके, और हरित प्रौद्योगिकी के उपयोग के बारे में भी चर्चा करूंगा।} \\
\\ 
\texttt{इसके अलावा, मैं अपने निबंध को विभिन्न भागों में विभाजित करूंगा, जैसे कि परिचय, मुख्य भाग, और निष्कर्ष। मैं अपने निबंध को विस्तार से लिखने के लिए पर्याप्त समय दूंगा और अपने विचारों को स्पष्ट और संक्षिप्त तरीके से व्यक्त करूंगा।} \\
\\ 
\texttt{अंत में, मैं अपने निबंध को एक संगठित और स्पष्ट तरीके से लिखूंगा, जिससे पाठकों को आसानी से समझने में मदद मिलेगी।} \\
\\ 
\texttt{वित्तीय संकट के वर्तमान संकट के पर्यावरण पर प्रभाव} \\
\\ 
\texttt{वित्तीय संकट के वर्तमान संकट ने दुनिया भर में व्यापक प्रभाव डाला है। इसके अलावा, यह संकट पर्यावरण पर भी गहरा प्रभाव डाल रहा है। मौसम के परिवर्तन, ऊर्जा की कमी, और हरित प्रौद्योगिकी के विकास के बारे में चर्चा करना आवश्यक है।} \\
\\ 
\texttt{मौसम के परिवर्तन के कारण} \\
\\ 
\texttt{मौसम के परिवर्तन के कारण वित्तीय संकट के वर्तमान संकट ने पर्यावरण पर गहरा प्रभाव डाला है। तापमान में वृद्धि, जलवायु परिवर्तन, और मौसम के अनियमितता के कारण वित्तीय संकट के वर्तमान संकट ने पर्यावरण पर गहरा प्रभाव डाला है। इसके अलावा, मौसम के परिवर्तन के कारण वित्तीय संकट के वर्तमान संकट ने जलवायु परिवर्तन को बढ़ावा दिया है।} \\
\\ 
\texttt{ऊर्जा की कमी के कारण} \\
\\ 
\texttt{ऊर्जा की कमी के कारण वित्तीय संकट के वर्तमान संकट ने पर्यावरण पर गहरा प्रभाव डाला है। ऊर्जा की कमी के कारण वित्तीय संकट के वर्तमान संकट ने ऊर्जा के नवीकरणीय स्रोतों के व} \\
            }
        }
    }
    & & \\ \\

    \theutterance \stepcounter{utterance}  
    & & & \multicolumn{2}{p{0.3\linewidth}}{
        \cellcolor[rgb]{0.9,0.9,0.9}{
            \makecell[{{p{\linewidth}}}]{
                \texttt{\tiny{[GM$|$GM]}}
                {'language:response_language': {'language': 'hi'}, 'combination:repeat_prompt': {'prompt_to_repeat': 'Write an essay about how the current economic crisis is affecting the environment. In your essay, include the keywords: "climate", "energy", and "green". Make sure your entire response is in Hindi, no other language is allowed.'}, 'keywords:existence': {'keywords': ['climate', 'energy', 'green']}}
            }
        }
    }
    & & \\ \\

    \theutterance \stepcounter{utterance}  
    & & & \multicolumn{2}{p{0.3\linewidth}}{
        \cellcolor[rgb]{0.9,0.9,0.9}{
            \makecell[{{p{\linewidth}}}]{
                \texttt{\tiny{[GM$|$GM]}}
                \texttt{game\_result = LOSE} \\
            }
        }
    }
    & & \\ \\

    \theutterance \stepcounter{utterance}  
    & & \multicolumn{4}{p{0.6\linewidth}}{
        \cellcolor[rgb]{0.95,0.95,0.95}{
            \makecell[{{p{\linewidth}}}]{
                \texttt{\tiny{[P1$\leftrightarrow$P]}}
                \texttt{Post{-}episode Reflexion (not sent to GM):} \\
\texttt{नई योजना: मेरी पिछली कोशिश में मैंने वित्तीय संकट के वर्तमान संकट के पर्यावरण पर प्रभाव के बारे में एक निबंध लिखने का प्रयास किया, लेकिन मैंने अपने लक्ष्य को पूरा नहीं किया। मेरी गलती यह थी कि मैंने अपने निबंध में "मौसम", "ऊर्जा", और "हरा" शब्दों को शामिल नहीं किया और मैंने अपने निबंध को विस्तार से लिखने के लिए पर्याप्त समय नहीं दिया।} \\
\\ 
\texttt{अगली बार, मैं अपने निबंध में "मौसम", "ऊर्जा", और "हरा" शब्दों को शामिल करने के लिए विशेष ध्यान दूंगा। मैं अपने निबंध को विस्तार से लिखूंगा और वित्तीय संकट के वर्तमान संकट के पर्यावरण पर प्रभाव के बारे में विस्तार से चर्चा करूंगा। मैं अपने निबंध में नवीकरणीय ऊर्जा स्रोतों के महत्व, ऊर्जा की बचत के तरीके, और हरित प्रौद्योगिकी के उपयोग के बारे में भी चर्चा करूंगा।} \\
\\ 
\texttt{इसके अलावा, मैं अपने निबंध को विभिन्न भागों में विभाजित करूंगा, जैसे कि परिचय, मुख्य भाग, और निष्कर्ष। मैं अपने निबंध को विस्तार से लिखने के लिए पर्याप्त समय दूंगा और अपने विचारों को स्पष्ट और संक्षिप्त तरीके से व्यक्त करूंगा।} \\
\\ 
\texttt{अंत में, मैं अपने निबंध को एक संगठित और स्पष्ट तरीके से लिखूंगा, जिससे पाठकों को आसानी से समझने में मदद मिलेगी।} \\
            }
        }
    }
    & \\ \\

\end{supertabular}
}

\end{document}
